%------------------------------------------------------------------------------%
\documentclass[fleqn,utf8,aspectratio=169,12pt]{beamer}

\usepackage{calculo_varias_variaveis-1}

\title{Funções de Várias Variáveis}

%------------------------------------------------------------------------------%
\begin{document}

\addtitle

%------------------------------------------------------------------------------%
\section{Funções}

%------------------------------------------------------------------------------%
\begin{frame}{Funções}
  \[
    f\colon A \to B
  \]
  \pause
  \vspace{-\baselineskip}
  \begin{itemize}[<+->]
    \item $A$ \emph{domínio}
    \item $B$ \emph{contradomínio}
    \item $f$ \emph{relação} que associa elementos do domínio a elementos do contradomínio
    \item Cada elemento do domínio deve estar associado a \emph{um único} elemento do contradomínio
    \item A \emph{imagem} é um subconjunto de $B$ composto dos \emph{elementos que são imagem} de algum elemento de $A$
  \end{itemize}
\end{frame}

%------------------------------------------------------------------------------%
\begin{frame}{Domínio}
  \begin{itemize}[<+->]
    \item Precisa ser dado a priori
    \item Geralmente é determinado pela aplicação
    \item Se nada for dito, assumimos o domínio ``natural'', \\
          ie, todos os valores onde podemos efetuar as contas
  \end{itemize}
\end{frame}

%------------------------------------------------------------------------------%
\begin{frame}{Imagem}
  \begin{itemize}[<+->]
    \item Subconjunto do contradomínio
    \item O contradomínio é definido a priori
    \item Pode ser difícil identificar a imagem de uma função
    \item Dada uma função \(f\colon D \to C \) \\
          $y$ pertence a imagem de $f$,
          se existe $x\in D$ tal que
          \[y=f(x)\]
  \end{itemize}
\end{frame}

%------------------------------------------------------------------------------%
\begin{frame}{Nomenclatura -- Curiosidade}
  Função real
  \[
    f\colon D \subset \R \to \R
  \]
  \pause
  Função de várias variáveis
  \[
    f\colon D \subset \R^n \to \R
  \]
  \pause
  Função vetorial ou paramétrica
  \[
    f\colon D \subset \R \to \R^n
  \]
  \pause
  Campo vetorial
  \[
    f\colon D \subset \R^n \to \R^m
  \]
\end{frame}

%------------------------------------------------------------------------------%
\section{Funções de Várias Variáveis}

%------------------------------------------------------------------------------%
\begin{frame}{Funções de Várias Variáveis}
  \[
    f\colon D \subset \R^n \to \R
  \]
  \pause
  \[
    f(x, y) = x + y
  \]
  \pause
  \[
    f(\vec{r}) = \sqrt{\sum_{k=1}^n r_k^2\,}
  \]
\end{frame}

%------------------------------------------------------------------------------%
\begin{frame}{Funções de Várias Variáveis}
  Suponha que $D$ seja um conjunto  de $n$-uplas de números reais
  \((x_1, x_2, \dots, x_n)\)

  \pause
  Uma \emph{função} a valores reais $f$ em $D$ é uma regra que associa
  \emph{um único número real}
  \[
    w = f(x_1, x_2, \dots, x_n)
  \]
  a cada elemento em $D$

  \pause
  O conjunto $D$ é o \emph{domínio} da função

  \pause
  O conjunto de valores de $w$ assumidos por $f$ é a \emph{imagem} da função

  \pause
  O símbolo $w$ é a variável \emph{dependente} de $f$, e dizemos que $f$
  é função de $n$ variáveis \emph{independentes} $x_1$ a $x_n$

  \pause
  Também chamamos os $x_j$ de variáveis de \emph{entrada} da função e
  denominamos $w$ a variável de \emph{saída} da função
\end{frame}

%------------------------------------------------------------------------------%
%\framedgraphic{Funções de Várias Variáveis}{definicao_funcao}{}

%------------------------------------------------------------------------------%
%\begin{frame}{Domínio e Imagem}
%
%  Determinamos o \emph{domínio} eliminando pontos que levam a divisões por zero ou
%  valores complexos.
%
%  \pause
%  \vspace{\baselineskip}
%
%  A \emph{imagem} consiste on conjunto de valores de saída para a variável dependente
%\end{frame}

%------------------------------------------------------------------------------%
\begin{frame}{Intervalos}
  Subconjuntos contínuos de $\R$
  \vspace{\baselineskip}

  \pause
  Intervalo fechado
  \[
    [0, 1] = \left\{ x\in\R\st 0 \leq x \leq 1 \right\}
  \]
  \pause
  Intervalo aberto
  \[
    (-2, 2) = \left\{ x\in\R\st -2 < x < 2 \right\}
  \]
  \pause
  Aberto a esquerda e fechado a direita
  \[
    (1, 5] = \left\{ x\in\R\st 1 < x \leq 5 \right\}
  \]
\end{frame}

%------------------------------------------------------------------------------%
\begin{frame}{Interior e Fronteira de um Conjunto}
  Um ponto $(x_0, y_0)$ em uma região (conjunto) $R$ no plano $xy$ é um
  \emph{ponto interior} de $R$ se é o centro de um disco de raio positivo que está
  inteiramente em $R$

  \pause
  \vspace{\baselineskip}
  Um ponto $(x_0, y_0)$ é um \emph{ponto de fronteira} de $R$ se todo disco
  centrado em $(x_0, y_0)$ contém ao mesmo tempo pontos que estão em $R$ e fora de $R$

  \pause
  \vspace{\baselineskip}
  O ponto de fronteira não precisa pertencer a $R$
\end{frame}

%------------------------------------------------------------------------------%
\begin{frame}{Interior e Fronteira de um Conjunto}
  Os pontos interiores de uma região, como um conjunto, \\
  formam o \emph{interior} da região

  \pause
  \vspace{\baselineskip}
  Os pontos de fronteira da região formam sua \emph{fronteira}

  \pause
  \vspace{\baselineskip}
  Uma região é \emph{aberta} se consiste apenas de pontos interiores

  \pause
  \vspace{\baselineskip}
  Uma região é \emph{fechada} se contem todos os seus pontos de fronteira
\end{frame}

%------------------------------------------------------------------------------%
%\framedgraphic{Regiões no Plano}{definicao_regioes}{}
\framedgraphic{Regiões no Plano}{pontos_regioes}{}

%------------------------------------------------------------------------------%
\begin{frame}{Regiões Limitadas}
  Uma região do plano é limitada se está dentro de um disco de raio fixo

  \pause
  \vspace{\baselineskip}

  Caso contrário ela é ilimitada
\end{frame}

% 01
%------------------------------------------------------------------------------%
\begin{question}
  Descreva o domínio da função \(f(x, y) = \sqrt{y-x^2}\)

  \pause
  \vspace{\baselineskip}
  \begin{enumerate}[<+->]
    \item Qual a condição que os pontos do domínio deve satisfazer?
    \item Que região é essa? Represente-a graficamente.
    \item Qual o interior da região?
    \item Qual a fronteira da região?
    \item A região é aberta ou fechada?
    \item A região é limitada?
  \end{enumerate}
\end{question}

% 01
%------------------------------------------------------------------------------%
\begin{resolution}{Solução}
  \onslide<+->{Condição que deve ser satisfeita para um ponto estar no domínio}
  \begin{align*}
    \onslide<+->{y-x^2 & \geq 0 \\}
    \onslide<+->{y & \geq x^2}
  \end{align*}

  \onslide<+->{
  Domínio
  \[
    D_f = \left\{ (x, y)\in\R^2 \st y \geq x^2\right\}
  \]}
\end{resolution}

%------------------------------------------------------------------------------%
\begin{resolution}{Representação Gráfica}
  \centering
  \includegraphics[height=0.8\textheight]{dominio_funcao}
\end{resolution}

%------------------------------------------------------------------------------%
\begin{resolution}{Solução}
  \begin{enumerate}
    \item \onslide<+->{Interior da região} \\
    \onslide<+->{
    \(
      y > x^2
    \)}
    \item \onslide<+->{Fronteira da região} \\
    \onslide<+->{
    \(
      y = x^2
    \)}
    \item \onslide<+->A região é aberta ou fechada? \\
    \onslide<+->{Fechada,}
    \onslide<+->{pois \(y = x^2\) pertence ao domínio}
    \item \onslide<+->{A região é limitada?} \\
    \onslide<+->{Ilimitada,}
    \onslide<+->{pois dado qualquer disco sempre teremos um ponto fora dele}
  \end{enumerate}
\end{resolution}

% 02
%------------------------------------------------------------------------------%
\begin{question}
  Encontre e esboce o domínio da função
  \[
    f(x, y) = \ln\left(x^2 + y^2 - 4\right)
  \]
\end{question}

%------------------------------------------------------------------------------%
\begin{resolution}{Solução}
  \onslide<+->{Condição}
  \begin{align*}
    \onslide<+->{x^2 + y^2 - 4 & > 0 \\}
    \onslide<+->{x^2 + y^2 & > 4}
  \end{align*}

  \onslide<+->{Pontos exteriores ao circulo centrado na origem e raio 2}

  \onslide<+->{
  Domínio
  \[
    D_f = \left\{ (x, y)\in\R^2 \st x^2 + y^2 > 4\right\}
  \]}
\end{resolution}

% 03
%------------------------------------------------------------------------------%
\begin{question}
  Encontre e esboce o domínio da função
  \(
    f(x, y) = \frac{\sqrt{x^2-4\,}}{\;\ln(y-1)\;}
  \)
\end{question}

%------------------------------------------------------------------------------%
\begin{resolution}{Solução}
  Condições para avaliar a função
  \(f(x, y) = \frac{\sqrt{x^2-4\,}}{\;\ln(y-1)\;}\)

  \pause
  \begin{enumerate}[<+->]
    \item \(x^2-4 \geq 0\)
    \item \(\ln(y-1) \neq 0\)
    \item \(y-1 > 0\)
  \end{enumerate}
\end{resolution}

%------------------------------------------------------------------------------%
\begin{resolution}{Solução}
\begin{columns}[t]
\column{0.3\linewidth}
  \onslide<+->{Condição 1}
  \begin{align*}
    \onslide<+->{x^2-4     & \geq 0                       \\}
    \onslide<+->{x^2       & \geq 4                       \\}
    \onslide<+->{\abs{x}   & \geq 2                       \\}
    \onslide<+->{x \leq -2 & \quad\text{ou}\quad x \geq 2}
  \end{align*}
\column{0.3\linewidth}
  \onslide<+->{Condição 2}
  \begin{align*}
    \onslide<+->{\ln(y-1) & \neq 0 \\}
    \onslide<+->{y-1      & \neq 1 \\}
    \onslide<+->{y        & \neq 2}
  \end{align*}
\column{0.3\linewidth}
  \onslide<+->{Condição 3}
  \begin{align*}
    \onslide<+->{y-1 & > 0 \\}
    \onslide<+->{y   & > 1}
  \end{align*}
\end{columns}
\end{resolution}

%------------------------------------------------------------------------------%
\begin{resolution}{Solução}
  Domínio
  \[
    D_f = \left\{ (x, y)\in\R^2 \st
      \abs{x} \geq 2 \;\text{ e }\;
      y > 1          \;\text{ e }\;
      y \neq 2
    \right\}
  \]
\end{resolution}

%------------------------------------------------------------------------------%
\section{Lista Mínima}

%------------------------------------------------------------------------------%
\begin{frame}{Lista Mínima}
  Cálculo Vol. 2 do Thomas 12\textsuperscript{a} ed. -- Seção 14.1
  \begin{enumerate}
    \item Estudar o texto da seção
    \item Resolver os exercícios: 1-2, 5-9
  \end{enumerate}

  \vfill
  Atenção: A prova é baseada no livro, não nas apresentações
\end{frame}

%------------------------------------------------------------------------------%
\end{document}
%------------------------------------------------------------------------------%
